\section{Introduction}\label{sec:intro}

In fraud detection, medical screening, intrusion detection and churn prediction, the class of interest is usually
the rarest. A classifier that minimises overall error can then reach high accuracy while missing most minority cases,
which keeps class imbalance a central problem in applied machine learning~\cite{he2009learning,krawczyk2016learning,
branco2016survey,fernandez2018book}. The standard remedies are \emph{data-level} resampling, which generates synthetic
minority samples or removes majority ones~\cite{chawla2002smote,he2008adasyn,batista2004study}, and ensembles that
reduce the variance and bias of minority predictions~\cite{galar2012review,chawla2003smoteboost,liu2009exploratory}.
Hybrid pipelines that resample the data and then fit an ensemble dominate recent applied work, particularly on
credit-card fraud~\cite{plan_p7_comparative,plan_p9_fraud,plan_p10_stacking}.

Our literature review (Section~\ref{sec:related}) found three recurring gaps. First, one resampler is chosen by hand
and applied \emph{uniformly} to the whole minority class, although the best resampler depends on the
dataset~\cite{plan_p8_hybrid,kovacs2019empirical} and plausibly on the region: small disjuncts~\cite{jo2004class},
borderline points~\cite{han2005borderline} and overlap zones~\cite{plan_p2_dcshs} each call for different treatment.
Second, ensemble combiners are static and trust every learner equally~\cite{plan_p5_survey}. Third, multiclass
imbalance is under-studied~\cite{plan_p3_multiclass,plan_p6_smoteif}.

We designed \ars{} to address all three. It clusters each minority class, scores the clusters for overlap and
density, picks a resampler per cluster, and combines Random Forest~\cite{breiman2001random},
XGBoost~\cite{chen2016xgboost} and LightGBM~\cite{ke2017lightgbm} with confidence weights. It supports binary and
multiclass problems and ships as a reproducible MLOps system.

Evaluating it fairly proved harder than building it. We found, and removed from our own pipeline, five pitfalls that
each manufacture a spurious improvement: unequal oversampling budgets, confidence weights cancelled by a meta-learner,
rank-forced selection rules, fixed decision thresholds, and duplicate records that straddle train and test splits.
With them removed, we asked whether adaptive resampler selection helps a modern tree ensemble at all. The answer is
largely negative, and we argue that this answer is the more useful result.

\para{Contributions}
\begin{enumerate}
  \item \textbf{Method:} \ars{}, with scale-free cluster scores, a global synthetic-sample budget, confidence-weighted
        averaging with a tuned threshold, and a multiclass extension (Section~\ref{sec:method}).
  \item \textbf{Protocol:} a leakage-free, budget-matched evaluation in which scaling, resampling, weighting and
        thresholds are fitted inside each fold (Section~\ref{sec:setup}).
  \item \textbf{Study:} nine variants on 28 datasets with repeated cross-validation, analysed with Friedman, Nemenyi,
        Holm-corrected Wilcoxon and corrected resampled $t$-tests, and with Bayesian tests using a region of
        practical equivalence~\cite{benavoli2017time} (Section~\ref{sec:results}).
  \item \textbf{Diagnosis:} why each component does or does not help, a validated-selection pilot, and a meta-analysis
        of when resampling helps (Section~\ref{sec:analysis}).
  \item \textbf{Artifact:} an open pipeline (DVC, MLflow~\cite{zaharia2018accelerating}, Prefect, Docker, CI,
        FastAPI) that records code version, packages and hardware with every result.
\end{enumerate}

In short, \ars{} beats aggressive cleaning (SMOTE-ENN) but is practically equivalent to plain SMOTE and to its
ablations, and it cannot be separated from no resampling. This agrees with growing evidence that strong learners
gain little from oversampling~\cite{elor2022smote,vandengoorbergh2022harm,vanhulse2007experimental}.
